\documentclass[letterpaper,10pt]{article}

% ================= PACKAGES =================
\usepackage[empty]{fullpage}
\usepackage{titlesec}
\usepackage{enumitem}
\usepackage[hidelinks]{hyperref}
\usepackage{fancyhdr}
\usepackage[none]{hyphenat}
\usepackage{fontawesome5}

\sloppy

% ================= PAGE SETUP =================
\pagestyle{fancy}
\fancyhf{}
\renewcommand{\headrulewidth}{0pt}

\addtolength{\oddsidemargin}{-0.6in}
\addtolength{\evensidemargin}{-0.6in}
\addtolength{\textwidth}{1.2in}
\addtolength{\topmargin}{-0.78in}
\addtolength{\textheight}{1.6in}

% ================= SPACING =================
\setlength{\parskip}{0pt}
\setlength{\itemsep}{0pt}
\setlength{\topsep}{0pt}
\setlength{\parsep}{0pt}

% ================= SECTION FORMAT =================
\titleformat{\section}
{\vspace{-8pt}\scshape\raggedright\large\bfseries}
{}{0em}{}[\titlerule\vspace{-5pt}]

% ================= CUSTOM COMMANDS =================
\newcommand{\projectheading}[3]{%
  \noindent\textbf{#1}\hfill{\small #2\quad#3}%
}

% ================= DOCUMENT =================
\begin{document}

% ================= HEADER =================
\begin{center}
{\Huge \scshape Divyansh Teja Edla} \\[3pt]
{\small Hyderabad, Telangana, India} \\[4pt]
{\small
\faPhone\ +91 9100324445 \quad
\faEnvelope\ \href{mailto:divyanshteja.edla@gmail.com}{divyanshteja.edla@gmail.com} \quad
\faGithub\ \href{https://github.com/DIVYANSH-TEJA-09}{DIVYANSH-TEJA-09} \quad
\faLinkedin\ \href{https://linkedin.com/in/divyansh-teja-edla-18557b28b}{divyansh-teja-edla}
}\\[2pt]
{\small
\faGlobe\ \href{https://divyansh-teja-edla.vercel.app}{\textbf{divyansh-teja-edla.vercel.app}} \quad
\faDatabase\ \href{https://huggingface.co/Divs0910}{huggingface.co/Divs0910}
}
\end{center}
\vspace{-6pt}

% ================= EDUCATION =================
\section{Education}
\textbf{B.E.\ in Computer Science and Engineering} \hfill \textbf{Expected: May 2026} \\
\textit{Matrusri Engineering College, Hyderabad} \hfill \textbf{CGPA: 8.89 / 10}

% ================= RESEARCH INTERESTS =================
\section{Research Interests}
Federated Learning, Low-Resource NLP (Indic Languages), Retrieval-Augmented Generation, Agentic AI, Explainable AI, Dataset Curation
\vspace{-2pt}

% ================= PUBLICATION =================
\section{Publication}
\textbf{Enhancing Federated Learning With Quantum-Inspired Particle Swarm Optimization: An IID MNIST Study} \\
\textit{Divyansh Teja Edla} and Dr.\ L.K.\ Indumathi \hfill NCICSET, 2025 \\
Proposed QPSO-based weight aggregation replacing FedAvg, achieving \textbf{97.6\% accuracy improvement} (81.51\% vs.\ 41.25\%).

% ================= RESEARCH EXPERIENCE =================
\section{Research Experience}
\textbf{Tech Lead — Viswam.AI} \hfill \textbf{June 2024 -- Present} \\
\textit{Low-Resource Language AI Initiative (Swecha $\times$ IIIT Hyderabad)}
\begin{itemize}[leftmargin=0.2in, itemsep=0pt, topsep=2pt, parsep=0pt]
\item Leading large-scale Telugu language dataset creation for LLM training and RAG systems.
\item Designed metadata extraction pipeline for a Telugu literary corpus (\textbf{758 books, 10,000+ stories}, 1947--2012) using Gemini 2.5 Flash API with structured JSON schemas for story-level indexing.
\item Built \textbf{Telugu OCR dataset}: \textbf{120K+ image-text pairs}, 10.7\,GB, published on \href{https://huggingface.co/datasets/Divs0910/telugu-ocr-dataset}{Hugging Face}.
\item Mentored \textbf{150+ students} contributing \textbf{1.5 lakh+ hours} of annotation and curation work.
\end{itemize}

% ================= PROJECTS =================
\section{Selected Projects}

\projectheading{Digi-Biz — Agentic Business Digitization}%
{\href{https://github.com/DIVYANSH-TEJA-09/business-digitization-agent}{\faGithub}}%
{\href{https://business-digitization-agent-two.vercel.app}{\faExternalLink*}}
\begin{itemize}[leftmargin=0.2in, itemsep=0pt, topsep=2pt, parsep=0pt]
\item \textbf{Problem:} 63M Indian small businesses lack digital presence; monolithic LLM calls fail on messy real-world documents.
\item Built \textbf{8-agent LangGraph pipeline} — file discovery $\rightarrow$ doc parsing $\rightarrow$ table extraction $\rightarrow$ media analysis $\rightarrow$ schema mapping $\rightarrow$ validation $\rightarrow$ profile output.
\item Achieved \textbf{95\% extraction completeness} (up from $\sim$60\% monolithic), processing 5+ input types in \textbf{$<$\,2 min}.
\item \textbf{Stack:} LangGraph, FastAPI, Groq/Llama-4-Scout, Streamlit, Docker.
\end{itemize}
\vspace{2pt}

\projectheading{FL-QPSO Brain Tumor AI Suite}%
{\href{https://github.com/DIVYANSH-TEJA-09/BrainTumor-FL-Pipeline}{\faGithub}}%
{\href{https://huggingface.co/spaces/Divs0910/Brain-Tumor-AI-Suite}{\faExternalLink*}}
\begin{itemize}[leftmargin=0.2in, itemsep=0pt, topsep=2pt, parsep=0pt]
\item \textbf{Problem:} FedAvg biases toward data-rich hospitals — rural clinics with fewer scans get the worst model.
\item Built \textbf{3D Attention U-Net} on BraTS 2021: \textbf{Dice 0.76} (WT), \textbf{0.85} (TC), \textbf{0.79} (ET) across 4 MRI modalities.
\item Replaced FedAvg with \textbf{QPSO-FL} that reduces cross-client Dice variance — optimizes for the floor, not the average. Benchmarked FedAvg vs FedProx vs QPSO-FL across 3 hospital nodes with Non-IID splits.
\item \textbf{Stack:} PyTorch, MONAI, QPSO, 3D Attention U-Net, ResNet-18, LSTM.
\end{itemize}
\vspace{2pt}

\projectheading{Telugu Agentic RAG System}%
{\href{https://github.com/DIVYANSH-TEJA-09/telugu-agentic-rag}{\faGithub}}%
{\href{https://telugu-story-generator-agent.streamlit.app}{\faExternalLink*}}
\begin{itemize}[leftmargin=0.2in, itemsep=0pt, topsep=2pt, parsep=0pt]
\item \textbf{Problem:} Telugu AI produces ``translated English'' — grammatically passable but culturally hollow.
\item Semantic search with GTE Multilingual + Qdrant over \textbf{10,000+ stories}: \textbf{92\% Hit Rate @\,1}, \textbf{99\% @\,5}, \textbf{MRR $\sim$0.95}.
\item Designed \textbf{Validator Agent} that rejects drafts with passive voice, lazy descriptions, or cultural inaccuracies — quality enforced architecturally, not by prompting.
\item \textbf{Stack:} Multi-agent, Qdrant, GTE Multilingual, Gemini API, Streamlit.
\end{itemize}

% ================= TECHNICAL SKILLS =================
\section{Technical Skills}
\vspace{1pt}
\begin{itemize}[leftmargin=0.2in, itemsep=0pt, topsep=0pt, parsep=0pt]
\item \textbf{Languages:} Python, C, R \quad \textbf{ML/DL:} PyTorch, TensorFlow, MONAI, XGBoost, SHAP, Scikit-learn
\item \textbf{NLP \& LLMs:} LangChain, LangGraph, HuggingFace, RAG, Prompt Engineering, Agentic Workflows
\item \textbf{Infra:} Qdrant, Chroma, FAISS, Docker, FastAPI, Streamlit, OpenAI/Gemini/Groq APIs, Ollama
\end{itemize}

% ================= LEADERSHIP =================
\section{Leadership}
\textbf{Founder \& President — DevCatalyst} \hfill \textbf{Aug 2025 -- Present} \\
Built 40-member technical community from scratch. Organized \textbf{8+ events} reaching \textbf{1,500+ students} — workshops on Agentic AI, Web3, Cloud (AWS), and Smart India Hackathon internal round.

% ================= HONORS =================
\section{Honors}
\textbf{Google Gemini Student Ambassador} (2025--26) \quad $\mid$ \quad \textbf{English Proficiency: C1 Level}

\end{document}
